% !TEX root = ../main.tex
\subsubsection*{Tytuł pracy}
\Title

\subsubsection*{Streszczenie}
W ramach pracy magisterskiej opracowano i poddano ewaluacji system automatycznej detekcji nadużyć w komunikacji tekstowej użytkowników platformy C2C. Analiza obejmuje identyfikację prób obejścia systemu płatności, wyłudzeń finansowych oraz treści toksycznych pojawiających się w rozmowach użytkowników.

W pracy porównano podejścia regułowe oraz rozwiązania wykorzystujące modele językowe, w tym model HerBERT. Oceniono skuteczność poszczególnych metod oraz możliwość ich zastosowania w rzeczywistym środowisku platform handlu internetowego.

\subsubsection*{Słowa kluczowe}
Uczenie maszynowe, przetwarzanie języka naturalnego, detekcja nadużyć, platforma C2C, HerBERT, klasyfikacja tekstu

\subsubsection*{Thesis title}
\begin{otherlanguage}{british}
\TitleAlt
\end{otherlanguage}

\subsubsection*{Abstract}
\begin{otherlanguage}{british}
This master's thesis presents the development and evaluation of a system for the automatic detection of abuse in text-based communication between users of C2C platforms. The analysis focuses on identifying payment system evasion attempts, financial fraud, and toxic content occurring in user conversations.

The study compares rule-based approaches with solutions based on language models, including HerBERT. The effectiveness of the proposed methods and their applicability in real-world e-commerce environments are evaluated.
\end{otherlanguage}
\subsubsection*{Key words}
\begin{otherlanguage}{british}
Machine learning, natural language processing, abuse detection, C2C platform, HerBERT, text classification
\end{otherlanguage}
